\chapter{Introduction}

\section{Problem statement}

The development of interactive computing and container orchestration has progressively changed how computational environments are provided to users. In 2014, the Jupyter Project branched off from the IPython project as a project aimed at supporting interactive data science and scientific computing across all programming languages. Around the same time that Jupyter was emerging, the open source project Kubernetes was introduced as a system to manage containerized workloads. Kubernetes received its first public commit in June 2014 and reached 1.0 in July 2015. JupyterHub, a multi-user version of the Jupyter notebook, was released as a preview in March 2015. The combination of JupyterHub and container orchestration technologies such as Kubernetes has since provided a practical foundation for delivering isolated and centrally managed notebook environments to groups of users. Today, JupyterHub explicitly supports Kubernetes-based deployments and is designed to provide users with computational environments and resources while allowing administrators to manage the underlying infrastructure.

Interactive computational environments have become an important part of modern data analysis, scientific computing, education, and machine learning. Jupyter notebooks are particularly useful in these settings because they allow users to combine executable code, computational results, and explanatory text within a single interactive environment. In multi-user deployments, however, providing notebook environments is not only a matter of running the notebook application itself. The platform must also manage users, compute resources, storage, software environments, and isolation between workloads. JupyterHub provides a common platform for deploying and managing multi-user Jupyter environments, while Kubernetes can provide the underlying infrastructure for scalable and isolated execution. In a typical JupyterHub deployment on Kubernetes, KubeSpawner creates a Kubernetes Pod for each user's notebook server and can expose a set of administrator-defined resource profiles and container images. Such profiles allow administrators to control the available CPU, memory, accelerators, and software environments while preventing users from directly requesting arbitrary infrastructure configurations. This administrator-controlled approach provides an important operational boundary, but it also introduces a usability problem. Users generally describe their requirements in terms of the work they want to perform rather than in terms of infrastructure configuration. For example, a user may know that they need to train a machine-learning model on a large dataset, process images, or perform data analysis with a particular framework. The platform, however, may require the user to select an appropriate resource profile and notebook image from a predefined catalog. Consequently, users may need to understand technical differences between available environments before they can start their actual work. This creates a gap between workload intent and infrastructure configuration. A user understands the intended computation, while the platform requires a concrete profile and image selection. Bridging this gap is particularly relevant when a catalog contains multiple environments with different resource capacities and software capabilities.

One way to present users with the options of administrator-defined solution profiles is to create a simple list that users can select from, one at a time. But that way users will have to interpret the solution catalog for themselves. A simple name for a profile does not necessarily imply that the solution will be suitable for a given workload. And relying on familiar software names to pick environments does not necessarily guarantee that the chosen environment meets the user’s explicit resource or capability requirements. Automating this decision is not straightforward either. A purely rule-based approach can provide deterministic and easily explainable decisions, but its ability to handle the variety of natural-language workload descriptions may be limited. Conversely, semantic retrieval or language-model-based approaches can improve the ability to interpret natural-language requests, but semantic relevance alone cannot guarantee that a candidate satisfies mandatory requirements or conforms to administrator-defined policies. The problem addressed by this thesis is therefore the design of a context-aware, policy-bounded mechanism for recommending an appropriate JupyterHub resource profile and notebook image from an administrator-controlled catalog based on a user's workload intent. The problem is deliberately narrower than general cloud resource optimization. This system does not attempt to predict exactly how much resources a given workload will use to run, replace the scheduling provided by Kubernetes, or try to optimize the whole cluster. What this system tries to do is decide on which of the pre-authorized environments by an administrator a user can select from their descriptions of intended workloads, and lists all relevant environments that satisfy all given requirements and constraints of the platform.

\section{Contributions and thesis overview}

The main contributions of this thesis are twofold. First, the thesis develops a context-aware resource profile selection approach that connects natural-language workload descriptions with a finite catalog of administrator-approved JupyterHub environments. The approach considers workload intent, explicit resource requirements, and the available environment catalog as relevant context during the selection process. This allows users to describe their computational needs at the workload level while keeping the actual resource environments within the set of configurations defined and controlled by the platform administrator. Second, the thesis proposes a multi-dimensional evaluation framework for assessing the resource profile selection approach. Rather than relying on a single recommendation metric, the evaluation considers several aspects of system behavior, including recommendation quality, robustness to variations in workload descriptions, resource-related behavior, notebook-image functionality, and practical integration with the JupyterHub environment. This broader evaluation is intended to provide a more comprehensive view of the effectiveness and practical applicability of the proposed approach.

The remainder of the thesis is organized into four chapters. Chapter 2 presents the technical and conceptual foundations required for the proposed work. It introduces JupyterHub and Kubernetes-based notebook environments, resource profiles and container images, workload intent, information retrieval, semantic retrieval, and recommendation-related concepts. These foundations establish the terminology and technical background used throughout the subsequent chapters. Chapter 3 presents the design of the proposed resource profile selection approach in detail. It describes the overall system architecture, workload-intent representation, candidate retrieval and ranking process, constraint handling, administrator-controlled environment catalog, and integration with the JupyterHub pre-spawn workflow. The chapter also discusses the main design decisions and defines the scope and boundaries of the proposed system. Chapter 4 describes the implementation and experimental methodology used to assess the proposed approach. It presents the experimental setup, evaluation scenarios, baselines, and evaluation metrics, followed by the experimental results and their discussion. The chapter also examines the limitations and scope of the evidence obtained from the experiments. Chapter 5 summarizes the work and its main findings, discusses the limitations of the current approach, and presents directions for future development and further evaluation. Overall, this thesis investigates the use of workload intent as an interface for selecting appropriate notebook environments in JupyterHub on Kubernetes. The central aim is to reduce the gap between user-level descriptions of computational tasks and infrastructure-level environment selection while maintaining administrator control and adherence to explicit platform constraints.
